\documentclass[acmsmall,11pt,nonacm,natbib=false]{acmart}
\geometry{letterpaper,margin=1in,includehead=false,includefoot=false,heightrounded=false,footskip=24pt}
\settopmatter{printacmref=false,printccs=false,printfolios=true}
\setcopyright{none}
\acmDOI{}
\acmISBN{}
\usepackage{mathtools}
\newtheorem{theorem}{Theorem}
\newtheorem{lemma}[theorem]{Lemma}
\newtheorem{corollary}[theorem]{Corollary}
\newcommand{\Z}{\mathbb Z}
\newcommand{\F}{\mathbb F}
\newcommand{\E}{\mathbb E}
\newcommand{\Prob}{\mathbb P}
\newcommand{\polylog}{\operatorname{polylog}}
\newcommand{\ord}{\operatorname{ord}}
\newcommand{\ceil}[1]{\left\lceil #1\right\rceil}
\newcommand{\floor}[1]{\left\lfloor #1\right\rfloor}
\title{Truly Subquadratic 3SUM over Every Bounded-Exponent Abelian Group}
\author{Zijian Zeng}
\affiliation{\institution{UCSI University}\country{Malaysia}}
\email{1002266693@ucsiuniversity.edu.my}
\renewcommand{\shortauthors}{Zijian Zeng}
\begin{document}
\begin{abstract}
\normalsize
For every fixed integer $M\ge2$, we give a randomized truly subquadratic
algorithm for three-list zero-sum matching over every compactly represented
finite abelian group of exponent dividing $M$ and polynomial order. The
exponent saving depends on $M$ and is independent of the polynomial order
bound. The algorithm has a worst-case time cap, returns a verified witness,
and succeeds with probability at least $3/4$ on every fixed positive input.
It extends the known exponent-dividing-$p^2$ scope to all higher prime powers
and mixed-prime products. The main step combines simultaneous profiles of
original annihilator-coset sizes with an extremal divisor certificate. The
certificate controls both class branching and false positives of every
possible sum order, enabling one elementary-middle support product per
profile. We use the existing selected-entry thin-matrix algorithm of Alman
and Vassilevska Williams. Through Majdoddin's finite-representation reduction,
the result gives a truly subquadratic algorithm for every fixed
abelian-group-representable finite coordinate relation.
\end{abstract}
\maketitle
\clearpage

\section{The result and its mechanism}

Three-list 3SUM asks whether $x_1+x_2+x_3=0$ for some $x_i$ in three given
lists. The ambient additive group affects the available hashing: a uniform
homomorphism can preserve a vector of order two much more often than one of
large order. This makes an order-independent collision estimate unsuitable
for groups with several torsion layers. We resolve the bounded-exponent
question by adapting the hash to all order layers simultaneously.

The input group is given as $G_0=\prod_{j=1}^r\Z_{m_j}$, with each
$m_j\mid M$, and list elements are residue vectors. Trivial factors are
discarded. For every fixed $c>0$ we allow $|G_0|\le n^c$. Costs include group
arithmetic, construction of every auxiliary list and matrix, all queries,
random-bit generation, and recovery of original list indices. A charged
word has $O_{M,c}(\log n)$ bits; bit simulation contributes a polynomial
logarithmic factor.

\begin{theorem}\label{thm:main}
Fix an integer $M\ge2$, and put
\[
 C=\ceil{\log_2 M},\qquad
 \eta=\min\{1/18,1/(8C)\},\qquad
 \gamma=63/1000,\qquad \delta_M=\eta\gamma/2.
\]
There is a uniform classical randomized algorithm which, for every fixed
$c>0$, solves three-list 3SUM over every such $G_0$ of order at most $n^c$
and lists of size at most $n$ in worst-case time
$O_{M,c}(n^{2-\delta_M}\polylog n)$. It always rejects negative inputs,
returns only verified witnesses, and returns a witness on each fixed
positive input with probability at least $3/4$.
\end{theorem}

Majdoddin's Theorem 19.17 and Corollary 19.25 give deterministic algorithms
for the entire exponent-dividing-$p^2$ class, for each fixed prime $p$;
Open Problem 19.26 asks for the general bounded-exponent extension
\cite{Maj26}. Theorem~\ref{thm:main} supplies that extension with randomness
and a smaller explicit saving. For example, its scope includes arbitrary
products of $\Z_8$, and products with different prime divisors, with ranks
growing logarithmically with $n$.

The support-product construction and its selected-entry matrix engine come
from prior work. Alman and Vassilevska Williams supply the thin-product
algorithm \cite{AVW26}; Majdoddin's Lemmas 19.6--19.7 organize homomorphic
matching as support products and probes \cite{Maj26}. The contribution here
is the simultaneous choice of quotient and row counts. We profile original
coset occupancies for every divisor of $M$. A least feasible row count gives
lower bounds for all divisors and an upper bound witnessed by one divisor.
The former controls every false-positive order, while the latter controls
the number of hash classes. The two bounds make the inherited support
construction work for the full bounded-exponent class.

The resulting tractability principle also covers every fixed finite
coordinate relation represented by an equation in an abelian group
(Corollary~\ref{cor:relation}). Its finite-quotient and translation arguments
are Majdoddin's Lemmas 20.3(d) and 20.4(a); Remark 20.6 explicitly anticipates
this consequence of the general group algorithm \cite{Maj26}. Recent work
of Chan, Itakura, Ma and Yu gives another truly subquadratic approach to
integer 3SUM through subset convolution \cite{CIMY26}. Our algorithm uses
the thin-product route and its group-specific balancing argument.

\paragraph{Research methodology.}
OpenAI Codex materially contributed to the divisor-profile proof, literature
inspection, exact finite verification code, and manuscript drafting. A fresh
model context checked the proof and implemented separate finite checks.
These roles affect the mathematical and algorithmic sections throughout;
the final AI methodology statement gives further detail.

\section{Normalization and original-coset profiles}\label{sec:profiles}

Deduplicate each list, retaining an original index for each value. Embed
$\Z_m$ into $\Z_M$ by $x\mapsto (M/m)x$. This is an injective
homomorphism, so the embedded lists have exactly the same zero-sum
witnesses. We henceforth work in $G=\Z_M^r$. Since $r\le c\log_2 n$, the
enlarged order is at most $n^{c\log_2 M}$ and each evaluation or comparison
of a residue vector has polynomial logarithmic cost. The trivial group and
empty lists are handled directly.

If one list has size at most $\sqrt n$, enumerate its pairs with another
list and use a sorted lookup in the third. The cost is
$O_{M,c}(n^{3/2}\polylog n)$. Assume below that all three original normalized
lists $S_i$ have size greater than $\sqrt n$. We also assume $n\ge16$ and
$R\ge M$ for the parameter $R$ defined in Section~\ref{sec:certificate}.
For fixed $M$ the excluded values of $n$ are bounded, so exhaustive search
on them is absorbed into the claimed bound.

Let $\mathcal D=\{d\ge2:d\mid M\}$ and $L=|\mathcal D|$. For $d\in\mathcal D$
write
\[
 H_d=G[d]=\{x:dx=0\}=(M/d)G.
\]
Reduction of coordinates modulo $M/d$ represents the quotient by $H_d$.
For each $x\in S_i$, let $c_{i,d}(x)$ be the number of elements of the
\emph{original} $S_i$ in $x+H_d$, and let $s_{i,d}(x)$ be the largest power
of two not exceeding $c_{i,d}(x)$. Partition $S_i$ by the full vector
$(s_{i,d}(x))_{d\in\mathcal D}$. Sorting computes all original coset counts
and these profiles in $O_{M,c}(n\polylog n)$ time. There are at most
$(1+\floor{\log_2 n})^L$ profiles per list.

Process every triple of nonempty profiles, obtaining a subinstance
$(U_1,U_2,U_3)$. Every original triple belongs to exactly one such subinstance.
Within it, $s_{i,d}$ is constant; define $s_d=\min_i s_{i,d}$. The following
two different uses of original occupancy are the reason the profile
restriction is safe.

\begin{lemma}\label{lem:profile}
For a fixed subinstance and each $d\in\mathcal D$:
\begin{enumerate}
\item the number $T_d$ of triples in $U_1\times U_2\times U_3$ with sum in
$H_d$ is less than $2s_d n^2$;
\item the number of original $H_d$-cosets represented in $U_i$ is at most
$n/s_{i,d}\le n/s_d$.
\end{enumerate}
Moreover $s_M>\sqrt n/2\ge n^{1/4}$.
\end{lemma}
\begin{proof}
Choose a list attaining $s_d$. Fixing the two other elements determines
one coset for its element. Any represented coset has fewer than $2s_d$
elements in its original list, and restriction cannot increase this count.
This proves the first assertion. For the second, each represented original
coset contains at least $s_{i,d}$ original elements; these cosets are
disjoint. Finally $H_M=G$, so each $s_{i,M}$ is the dyadic rounding of
$|S_i|>\sqrt n$. The last inequality follows from $n\ge16$.
\end{proof}

A restricted coset may contain only one element, even when its original
occupancy is large. Lemma~\ref{lem:profile} uses an upper bound on restricted
occupancy and a lower bound on original occupancy; it never needs a lower
bound on restricted occupancy.

\section{The divisor certificate and the hash}\label{sec:certificate}

For a fixed profile subinstance set
\begin{equation}\label{eq:parameters}
 h=\floor{\eta\log_2 n},\quad R=2^h,\quad
 \alpha=\tfrac12-\tfrac\gamma4=\tfrac{1937}{4000},\quad
 \sigma=\floor{R^\alpha}.
\end{equation}
Thus $n^\eta/2\le R\le n^\eta$ and
$R^\alpha/2\le\sigma\le R^\alpha\le\sqrt R$.
Choose the least nonnegative integer $\ell$ satisfying
\begin{equation}\label{eq:lower}
 d^\ell\sigma\ge s_dR\qquad(d\in\mathcal D).
\end{equation}
Because $s_M\ge n^{1/4}$ and $R/\sigma\ge1$, $\ell\ge1$.
Minimality supplies $e\in\mathcal D$ for which the inequality fails at
$\ell-1$, and hence
\begin{equation}\label{eq:upper}
 e^\ell< e s_e R/\sigma.
\end{equation}
This is the extremal divisor certificate: all divisors have the lower
bound~\eqref{eq:lower}, while one divisor has the complementary upper
bound~\eqref{eq:upper}. It is computed from the profile values in
$O_M(\log n)$ integer-power comparisons, and $\ell=O_M(\log n)$.

Choose a prime $p\mid e$, and put
\[
 q=\floor{\log_p R},\qquad D=p^q,\qquad k=\ell-q.
\]
Our large-$n$ assumptions imply $q\ge1$ and $R/M<D\le R$. Also
\[
 \ell\ge\tfrac14\log_M n,\qquad
 q\le\eta\log_2 n\le\tfrac18\log_M n,
\]
so $k\ge0$. All parameters are exact integers. For example, rational-power
rounding in~\eqref{eq:parameters} is implemented by integer-power comparisons
on $O_M(\log n)$-bit numbers.

Take independent uniform rows $u_1,\ldots,u_\ell\in\Z_M^r$. Define two
homomorphisms by
\begin{align}
 \pi(x)&=\bigl(x+H_e,
 (\langle u_j,x\rangle\bmod M)_{j\le k},
 (\langle u_j,x\rangle\bmod (M/p))_{k<j\le\ell}\bigr),\label{eq:pi}\\
 K(x)&=\bigl(x+H_e,
 (\langle u_j,x\rangle\bmod M)_{j\le\ell}\bigr).\label{eq:K}
\end{align}
The value of $K$ determines $\pi$. The quotient component is stored as a
residue vector; its entire group is never enumerated. Each evaluation has
$O_{M,c}(\log^2 n)$ cost. If $x\in\ker\pi$, the quotient component and
the first $k$ coordinates of $K(x)$ vanish, and each remaining coordinate
is a multiple of $M/p$. Consequently
\begin{equation}\label{eq:B}
 K(\ker\pi)\subseteq
 B=\{0\}\times\{0\}^k\times((M/p)\Z_M)^q\cong\F_p^q,
 \qquad |B|=D.
\end{equation}
Using this whole explicit subgroup is sufficient, even if the actual image
is smaller. Split its coordinates as $B=B_1\times B_2$ with dimensions
$\ceil{q/2}$ and $\floor{q/2}$.

\begin{lemma}\label{lem:blocks}
For every choice of rows, each $U_i$ has fewer than $M^2n/\sigma$ nonempty
$\pi$-classes. Cutting each such class into blocks of size at most $\sigma$
gives at most $(M^2+1)n/\sigma$ blocks.
\end{lemma}
\begin{proof}
On one $H_e$-coset a full row has at most $e$ values. A row reduced modulo
$M/p$ has at most $e/p$ values, since $p\mid e$. Thus at most
\[
 e^k(e/p)^q=e^\ell/D
 < e s_e R/(\sigma D)<M^2s_e/\sigma
\]
distinct $\pi$-labels occur in that coset. Multiply by at most $n/s_e$
represented original cosets, using Lemma~\ref{lem:profile}. The number of
blocks is at most the number of classes plus $|U_i|/\sigma$.
\end{proof}

\section{Controlling every sum order}\label{sec:orders}

\begin{lemma}\label{lem:survival}
If $v\in G$ has additive order $d$, a uniform row maps $v$ uniformly onto
the subgroup of $\Z_M$ of order $d$. For nonzero $v\in H_e$,
\begin{equation}\label{eq:survival}
 \Prob[K(v)=0]=d^{-\ell},\qquad
 \Prob[\pi(v)=0]=\gcd(d,p)^q d^{-\ell}\le Dd^{-\ell}.
\end{equation}
Vectors outside $H_e$ have survival probability zero under both maps.
\end{lemma}
\begin{proof}
The map $u\mapsto\langle u,v\rangle$ is a homomorphism. Its image is the
ideal generated by the coordinates of $v$ in $\Z_M$, of size $d$; all
fibers have equal size. Reduction modulo $M/p$ takes that image to one of
size $d/\gcd(d,p)$. Independence of the rows gives~\eqref{eq:survival}.
The common quotient component excludes vectors outside $H_e$.
\end{proof}

Let $Z_0$ count true zero-sum triples in the subinstance, let $Z_\pi$ count
triples whose sum is in $\ker\pi$, and let $F$ count triples of nonzero
sum in $\ker K$. Every nonzero element of $H_e$ has order $d>1$ dividing
$e$. The number of triples of exact sum-order $d$ is at most $T_d$, hence
less than $2s_dn^2$. Lemma~\ref{lem:survival} and~\eqref{eq:lower} give
\begin{align}
 \E F&\le\sum_{\substack{d\mid e\\d>1}}2s_dn^2d^{-\ell}
       \le 2L\sigma n^2/R,\label{eq:EF}\\
 \E(Z_\pi-Z_0)&\le D\E F\le2L\sigma n^2.\label{eq:EZ}
\end{align}
Markov's inequality and a union bound imply that with probability at least
$5/6$ both
\begin{equation}\label{eq:good}
 F\le24L\sigma n^2/R,\qquad Z_\pi-Z_0\le24L\sigma n^2
\end{equation}
hold. Dependence between these two random counts is harmless. Since lists
are sets, $Z_0\le n^2$. We call~\eqref{eq:good} the good event. The counts
in its definition are used only in the analysis. The algorithm observes
the listing and scanning caps below instead of evaluating these counts.

\section{Support products, probes, and recovery}\label{sec:support}

We spell out the support construction to account for all queries and
witness scans. The construction follows the framework of Majdoddin's
Lemmas 19.6--19.7 \cite{Maj26}. Form the blocks in lists 2 and 3 given by
Lemma~\ref{lem:blocks}, and choose a representative $z_b$ in every block.
A probe is $Q=(x_1;b_2,b_3)$ with $x_1\in U_1$ and
\[
 \pi(x_1)+\pi(z_{b_2})+\pi(z_{b_3})=0.
\]
For every pair $(x_1,b_2)$, look up the compatible label of list 3 and
enumerate its blocks. Sorted label indices implement this in polynomial
logarithmic time per pair and output.

\begin{lemma}\label{lem:probes}
Let $C_2$ be the number of nonempty $\pi$-classes in $U_2$. The set
$\mathcal P$ of probes satisfies
\begin{equation}\label{eq:probes}
 |\mathcal P|\le n C_2+2n^2/\sigma+Z_\pi/\sigma^2.
\end{equation}
On the good event, $|\mathcal P|\le A n^2/\sigma$ for
$A=M^2+3+24L$.
\end{lemma}
\begin{proof}
Fix $x_1$. For a compatible pair of nonempty classes of sizes $a,b$,
the block-pair count is at most
$(1+a/\sigma)(1+b/\sigma)$. Summing the constant term over all $x_1$
and labels gives at most $nC_2$. The two linear terms each sum to at most
$n^2/\sigma$: translation of labels is injective, so each class of either
list appears at most once for a fixed $x_1$. The bilinear term sums to
$Z_\pi/\sigma^2$. On the good event,
$Z_\pi\le n^2+24L\sigma n^2$; use Lemma~\ref{lem:blocks} and $\sigma\ge1$.
\end{proof}

During listing, abort this subinstance if the next probe would exceed
$\ceil{A n^2/\sigma}$. This cap is observable. The outer enumeration costs
$O_M(n^2/\sigma)$, plus polynomial logarithmic work per stored probe.
An aborted subinstance contributes no witness and processing continues.

For a nonaborted subinstance put
\[
 b_0=M^2+1,\qquad g_0=M(A+b_0+3),\qquad
 g'=\ceil{g_0n/\sigma},\qquad N=g'p^{\ceil{q/2}}.
\]
Left vertices are pairs $(b_2,\mu_1)$ with $\mu_1\in B_1$; right vertices
are $(b_3,\mu_2)$ with $\mu_2\in B_2$. Pad both sides to size $N$, and use
$B$ as a middle set. Lemma~\ref{lem:blocks} ensures sufficient room.
Construct $0/1$ matrices $X\in\Z^{N\times D}$ and
$Y\in\Z^{D\times N}$ by
\begin{align}
 X[(b_2,\mu_1),\lambda]=1
 &\iff \exists x_2\in b_2:\
 K(x_2)-K(z_{b_2})+(\mu_1,0)=\lambda,\label{eq:X}\\
 Y[\lambda,(b_3,\mu_2)]=1
 &\iff \exists x_3\in b_3:\
 -(K(x_3)-K(z_{b_3}))-(0,\mu_2)=\lambda.\label{eq:Y}
\end{align}
All differences in these definitions lie in $B$, because a block lies
within one $\pi$-class. For a probe $Q$ the element
\[
 w_Q=K(x_1)+K(z_{b_2})+K(z_{b_3})\in B
\]
has a unique decomposition $(\mu_1,\mu_2)$. Query the product at
$((b_2,\mu_1),(b_3,\mu_2))$. Deduplicate these matrix positions while
retaining the corresponding probe lists.

\begin{lemma}\label{lem:semantics}
The product entry for a probe is positive exactly when some pair
$x_2\in b_2$, $x_3\in b_3$ satisfies $K(x_1+x_2+x_3)=0$.
Every true witness has a positive probe, and at most $F$ positive probes
contain no true witness.
\end{lemma}
\begin{proof}
A positive integer product entry means that the two support equations
\eqref{eq:X}--\eqref{eq:Y} hold for one common $\lambda$. Equating them
gives
\[
 K(x_2)+K(x_3)-K(z_{b_2})-K(z_{b_3})+(\mu_1,\mu_2)=0.
\]
Substituting $w_Q$ yields precisely $K(x_1+x_2+x_3)=0$. Each step is
reversible. A true witness has zero $\pi$-sum and zero $K$-sum, and thus
gives a positive probe. A positive probe with no true witness contains
a nonzero-sum $K$-zero triple. Distinct probes contain disjoint triples,
since a triple determines $x_1$ and both blocks. Choosing one such triple
per failed probe bounds their number by $F$, even if probes share a matrix
position.
\end{proof}

Scan a positive probe by enumerating $x_2\in b_2$ and looking up
$-x_1-x_2$ in the sorted block $b_3$. This uses original residue vectors,
costs $O_{M,c}(\sigma\polylog n)$, and returns original indices when
successful. Scan at most $\ceil{24L\sigma n^2/R}+1$ positive probes,
stopping immediately on a verified witness. If all these scans fail,
abort the subinstance. On the good event, Lemma~\ref{lem:semantics}
ensures that a subinstance containing a witness finds one within this
cap. A final direct sum check verifies every returned witness.

\section{The legal matrix call and the total bound}\label{sec:time}

We use the following existing algorithmic theorem, including its stated
integer-input and selected-entry interface.

\begin{lemma}[Alman--Vassilevska Williams, Corollary 26 \cite{AVW26}]
\label{lem:thin}
Let $N\ge D^{18}$, and let $X\in\Z^{N\times D}$ and
$Y\in\Z^{D\times N}$ have entries of absolute value $N^{O(1)}$.
They can be preprocessed deterministically in $O(N^2/D^{0.063})$ time
and space, after which an entry of $XY$ can be queried in $O(D^{0.437})$
time. In particular, at most $N^2/\sqrt D$ prescribed entries can be
computed in $O(N^2/D^{0.063})$ time.
\end{lemma}

Here the entries of $X,Y$ are $0/1$. Since $\sigma\le\sqrt R$ and
$D>R/M$,
\[
 N\ge g_0n\sqrt D/\sigma\ge n.
\]
Consequently $D\le R\le n^\eta\le N^{1/18}$. Also
$D^{1+\gamma}\le N$, since $\eta\le1/18$ and $\gamma<1/2$.
The listing cap is at most $(A+1)n^2/\sigma$, and
\begin{equation}\label{eq:querybudget}
 \frac{N^2}{\sqrt D}\ge
 \frac{g_0^2n^2\sqrt D}{\sigma^2}\ge
 \frac{g_0^2}{\sqrt M}\frac{n^2}{\sigma}\ge
 (A+1)\frac{n^2}{\sigma}.
\end{equation}
Thus one call to Lemma~\ref{lem:thin} computes all required positions.

For the cost, $p^{\ceil{q/2}}\le\sqrt{pD}$ and
$g'\le(g_0+1)n/\sigma$. Since $D$ and $R$ differ by at most a fixed
factor, the matrix call costs
\begin{equation}\label{eq:matrixcost}
 O_M(N^2/D^\gamma)=O_M(n^2R^{1-\gamma}/\sigma^2)
 =O_M(n^2R^{-\gamma/2}).
\end{equation}
Initializing the two dense input matrices takes $O(ND)$ time, covered by
this bound because $D^{1+\gamma}\le N$. Set their nonzero entries by
enumerating each list element and all possible $\mu_i$. This costs
$O_{M,c}(n\sqrt{pD}\polylog n)$, also covered by~\eqref{eq:matrixcost}.
Sorting the capped probe positions and retaining their lists costs
$O_{M,c}(n^2\sigma^{-1}\polylog n)$. This fits because
$\alpha\ge\gamma/2$.

The capped scans, including the possible final successful scan, cost
\begin{equation}\label{eq:scancost}
 O_{M,c}\bigl((\sigma^2n^2/R+\sigma)\polylog n\bigr)
 =O_{M,c}(n^2R^{-\gamma/2}\polylog n).
\end{equation}
The equality in the dominant exponent follows from
$2\alpha-1=-\gamma/2$. Row evaluation, class sorting, label arithmetic,
parameter computation and recovery use $O_{M,c}(n\polylog n)$ additional
time. Matrix dimensions are polynomial in $n$, so word sizes and their
bit simulation add polynomial logarithmic overhead only.

For completeness, uniform coefficients in $\Z_M$ are generated from
$C$ fair bits by rejecting values at least $M$. An accepted coefficient
needs at most two attempts in expectation. For each profile subinstance,
cap the total number of attempts for its $r\ell$ coefficients at
$24r\ell$, aborting the subinstance if the cap is exceeded. This cap costs
$O_M(r\ell)$ time on every execution. Coupling with unlimited independent
uniform generation, Markov's inequality bounds the probability of such
an abort by $1/12$.

Every subinstance therefore has cost
$O_{M,c}(n^2R^{-\gamma/2}\polylog n)$ on every random path, whether or
not its listing or scans abort. There are at most
$(1+\floor{\log_2 n})^{3L}$ subinstances. As $M$ is fixed and
$R\ge n^\eta/2$, their total cost is the bound in
Theorem~\ref{thm:main}, with saving $\delta_M=\eta\gamma/2$ independent
of $c$.

To verify correctness, fix a witness of a positive input before any
randomness. Its profile subinstance is deterministic. Its ideal uniform
rows satisfy the two inequalities in~\eqref{eq:good} except with probability
at most $1/6$. The coefficient cap has probability at most $1/12$ of
interfering. On the complementary event the listing cap holds, and a
positive probe containing a true witness is reached within the scan cap.
Thus this subinstance finds a witness with probability at least $3/4$.
If the algorithm stops earlier it has already found a verified witness.
One may regard independent random tapes for all subinstances as sampled
in advance, so stopping early does not require a union bound over them.
Negative inputs always return absence because every success requires an
original zero-sum triple. This proves Theorem~\ref{thm:main}.

\section{Finite coordinate relations}\label{sec:relations}

Let $\Sigma$ be a fixed finite alphabet. A relation $P\subseteq\Sigma^3$
is abelian-group-representable if there are maps $\phi_i:\Sigma\to A$
and $t\in A$ into an abelian group with
\[
 (a,b,c)\in P\quad\Longleftrightarrow\quad
 \phi_1(a)+\phi_2(b)+\phi_3(c)=t.
\]
Given three lists in $\Sigma^d$, the coordinatewise problem asks for one
triple whose coordinate triples all belong to $P$.

\begin{corollary}\label{cor:relation}
For each fixed abelian-group-representable finite relation $P$, there is
$\delta_P>0$ such that, for every fixed $c>0$, its coordinatewise problem
on lists of size at most $n$ and $d\le c\log n$ has a randomized witness
algorithm of worst-case time $O_{P,c}(n^{2-\delta_P}\polylog n)$,
with the one-sided guarantee of Theorem~\ref{thm:main}.
\end{corollary}
\begin{proof}
We give the finite-quotient argument of Majdoddin's Lemma 20.3(d), followed
by the translation of Lemma 20.4(a) \cite{Maj26}. The empty relation is
immediate. Otherwise fix $s^0\in P$ and use the basis $e_{i,a}$ of
$\Z^{3|\Sigma|}$. Put
\[
 e(s)=e_{1,s_1}+e_{2,s_2}+e_{3,s_3},\qquad
 L_P=\operatorname{span}_{\Z}\{e(s)-e(s^0):s\in P\}.
\]
Any representation induces a homomorphism vanishing on $L_P$, so for
$s\notin P$ the vector $e(s)-e(s^0)$ is outside $L_P$. In a Smith normal
form basis, write $L_P=\bigoplus_{i\le\rho}d_i\Z e'_i$. There are only
finitely many forbidden vectors. Choose an integer $m\ge2$ divisible by
every $d_i$ and larger than the absolute values of all coordinates of
these vectors. A forbidden vector $y$ remains outside
$L_P+m\Z^{3|\Sigma|}$: either some $i\le\rho$ has $d_i\nmid y_i$, or
some $i>\rho$ has nonzero $y_i$, which cannot be divisible by $m$.
The finite group
\[
 A_P=\Z^{3|\Sigma|}/(L_P+m\Z^{3|\Sigma|})
\]
therefore represents exactly $P$, using $\phi_i(a)=[e_{i,a}]$ and
$t=[e(s^0)]$. This group and its exponent depend only on $P$.

For a string $x\in\Sigma^d$ define
$\Phi_i(x)=(\phi_i(x[j]))_{j=1}^d$. Translate the third list by
$-(t,\ldots,t)$ and leave the other two image lists unshifted. A resulting
zero sum in $A_P^d$ is equivalent, coordinate by coordinate, to a valid
original triple. The construction costs $O_P(nd)$; retain a preimage
index for every image value. These indices recover a valid original
witness even when image strings coincide. Since
$|A_P^d|\le n^{O_P(c)}$ and its exponent is fixed, apply
Theorem~\ref{thm:main} with $M=\max\{2,\exp(A_P)\}$ and
$\delta_P=\delta_M$. The exponent saving is independent of $c$.
\end{proof}

This supplies the sufficient tractability direction for every representable
relation. The argument leaves the complexity of nonrepresentable relations
as a separate question. The prior finite-quotient reduction accounts for
the breadth of the consequence; the group algorithm supplies the extension
that it requires.

\section{Verification and further questions}

The proof gives an explicit, bounded-time reduction to the prior
thin-product algorithm. Its finite checks supplement the algebra: the
separate checker tested 778 full-row and 1,382 reduced-row distributions,
4,622 profile subinstances, and 144 literal support instances containing
45,932 matrix entries and 13,244 probes. The tests included mixed moduli,
fragmented profiles and nonsurjective images. These are regression tests;
the bounds for all inputs follow from the lemmas above. The thin-product
algorithm is used as a cited theorem. A faster saving as $M$ grows and a
deterministic replacement for the row-selection step remain useful
directions.

\bibliographystyle{alpha}
\bibliography{references}

\section*{AI methodology statement}
OpenAI Codex generated and materially influenced the proof development,
literature comparisons, verification programs and English exposition in
this paper. Codex used primary-source text and web metadata for the cited
results. A separate fresh model context reconstructed the proof and wrote
a second exact checker; finite tests were executed locally. The manuscript
was drafted and compiled with Codex assistance. Proof verification consisted
of model-based mathematical reconstruction and local exact finite checks.
The sole listed author is responsible for the paper's content.
\end{document}
